\section{Transverse-Momentum Dependent PDFs}
\label{sec:tmd}
%------------------------------------------------------------------------------
The transverse-momentum-dependent (TMD) parton distribution functions (TMDPDFs) are a natural generalization
of the collinear PDFs to include both longitudinal and transverse momentum of partons.
They are in principle probability distributions $f_i(x,\vec{k}_\perp,\sigma)$ of finding a parton of given species $i$,
longitudinal and transverse momentum $(xP^+,\vec{k}_{\perp})$, and polarization $\sigma$
inside the hadron state. TMDPDFs are playing an increasingly important role in
understanding the partonic structure of hadrons and high-energy scattering.

The TMD parton densities were firstly introduced by Collins and Soper in 1980s~\cite{Collins:1981uk,Collins:1981va,Collins:1982wa,Collins:1984kg,Collins:1985ue,Bodwin:1984hc} to understand the Drell-Yan (DY) and $e^{+}e^{-}$ annihilation process, and generalized in~\cite{Ji:2004wu,Ji:2004xq} to semi-inclusive deep-inelastic scattering(SIDIS) process. The TMD factorization has been reanalyzed in the framework of SCET in which modes are made manifest by effective fields ~\cite{Bauer:2000yr,Bauer:2001yt,Manohar:2006nz,Becher:2010tm,GarciaEchevarria:2011rb,Echevarria:2012js,Chiu:2012ir}. Various TMD factorization formalisms finally converged to the standard one where a scheme-independent TMDPDF can be defined~\cite{Echevarria:2012js,Collins:2012uy,Collins:2017oxh}.

The TMD parton densities are important in understanding the experimental processes where the transverse momenta of final state particles are measured. For example, in DY pair and $W, Z$ production it is known that the differential cross section $d\sigma/dQ_T^2$ normally peaks at relatively small transverse momentum.  For $Q\sim 10 $ GeV, the peak is typically located at $Q_\perp\sim 1 $ GeV where nonperturbative effects are important~\cite{Collins:1984kg}. A good knowledge of TMD parton densities is therefore crucial for the determination of the
cross sections and precision test of perturbative QCD predictions.

Besides their importance in understanding the high-energy experimental data, the TMD parton densities are also important by themselves for their crucial role in describing hadron structures. With them, one can simultaneously study the fast-moving collinear physics through the longitudinal $x$-dependencies, and the nonperturbative effect from the transverse $\vec{k}_\perp$-dependencies. Moreover, the TMDPDFs are sensitive to effects such as soft radiations. Therefore, the physics in the presence of transverse degrees of freedom is rather rich. This is particularly true in studies of spin-dependent phenomena where one can define various TMDPDFs through Lorentz decompositions (see Sec. V.B). One example is the Sivers function for a transversely polarized proton, $\epsilon_{ij}k^i_\perp S^i_\perp f^{\perp}_{1T}(x,k_\perp)$, which is naive-time-reversal odd and is predicted to change sign between the DY and SIDIS processes~\cite{Collins:2011zzd}. Similar properties also exist in the Boer-Mulders function~\cite{Boer:1997nt} concerning a transversely-polarized parton
distribution in an unpolarized hadron. These two functions are related to the single transverse spin asymmetry. If we generalize the TMDPDFs to include the impact parameter dependence, we can further define the Wigner function, the parton orbital angular momentum distributions, etc~\cite{Belitsky:2003nz, Lorce:2011ni}. Therefore, the TMDPDFs allow for a more complete and refined 3D description (or tomography) of the hadron structure~\cite{Burkardt:2000za, Boer:2011fh}.
The 3D tomography of the proton is a major physical goal of the EIC program. The TMDPDFs are also important
in understanding small-$x$ physics~\cite{Kuraev:1977fs,Balitsky:1978ic,Balitsky:1995ub,Kovchegov:1999yj,Kovchegov:2012mbw}.

Our current knowledge on TMDPDFs mainly comes from fitting to the experimental data~\cite{Landry:1999an,Konychev:2005iy,Su:2014wpa,Echevarria:2014xaa,Kang:2015msa,Bacchetta:2017gcc,Scimemi:2017etj,Bertone:2019nxa,Scimemi:2019cmh,Bacchetta:2019sam}. This is, however, rather primitive due to the paucity of data. Although the future EIC will make up the gap and produce more data for TMD measurements, it is still important to develop first-principle methods for the determination of nonperturbative TMDPDFs, which can serve as a test or provide useful inputs to constrain the global fits.
LaMET provides a systematic way to extract TMDPDFs from the lattice calculations. Early studies~\cite{Ji:2014hxa,Ji:2018hvs,Ebert:2018gzl,Ebert:2019okf} have tried to construct a quasi-TMDPDF on the lattice, but its relation to the physical TMDPDF is expected to be nonperturbative due to complications in the soft function~\cite{Ebert:2019okf}. The recent works in~\cite{Ji:2019sxk,Ji:2019ewn} provide a formulation
to calculate the soft function so that a perturbative matching formula can be established between the quasi- and physical TMDPDFs, allowing for a complete determination of the latter from lattice QCD. In this section we review the application of LaMET to the nonperturbative TMDPDFs. The investigation is still in its early stage and a lot remains
to be explored, particularly in lattice calculations and matching.

In the first subsection we introduce the TMDPDFs and discuss the associated rapidity divergences. In the following subsections, we define the quasi-TMDPDFs or TMD momentum distributions in a proton of finite momentum,
and study their momentum RGEs and UV renormalization properties. In the process,
we introduce the off-light-cone soft functions. We then present the factorization of the quasi-TMDPDFs into the light-cone TMDPDFs and the off-light-cone soft function, where various one-loop results and the relevant RGEs are also given.
The properties of the off-light-cone soft function are discussed in the last subsection, where it is shown to be related to
the form factor of a pair of charged color sources, which paves the way for its calculation on a Euclidean lattice.

\subsection{Introduction to TMDPDFs and Rapidity Divergence}
\label{sec:lightconetmd}

As explained in \sec{lamet}, we can define various TMDPDFs by choosing different
gauge-links between the quark or gluon bilinears. The one relevant to high-energy phenomena is defined with
light-like Wilson lines. The links represent the propagation of high-energy color-charged particles, and are crucial in forming gauge-invariant nonlocal operators~\cite{Belitsky:2002sm}. As argued in previous sections, such
operators are the result of an EFT description (more explicitly so in SCET) arising from
taking the infinite-momentum limit of the proton. Thus, it is natural to expect that they require additional regularization and renormalization.

Let us take the non-singlet quark unpolarized TMDPDF as an example.
Without the field theoretic subtleties, the distribution is
\begin{align}\label{eq:beam}
	& f(x,\vec{k}_\perp) = \frac{1}{2P^+}\int\frac{d\lambda}{2\pi}\frac{d^2\vec{b}_\perp}{(2\pi)^2}
	e^{-i\lambda x+i\vec{k}_\perp \cdot \vec{b}_\perp}\\
	& \times \langle P| \bar \psi(\lambda n/2 +\vec{b}_\perp)\gamma^+{\cal W}_n(\lambda n/2+\vec{b}_\perp)\psi(-\lambda n/2)|P\rangle  \ ,\nn
\end{align}
where ${\cal W}_n(\lambda n+\vec{b}_\perp)$ is the staple-shaped gauge-link of the form
\begin{align}
	&{\cal W}_n(\xi)=W^{\dagger}_{n}(\xi)W_{\perp}W_{n}(-\xi \cdot p n) \label{eq:staplen} \ , \\
	&W_{n}(\xi)= {\cal P}\exp\left[-ig\int_{0}^{-\infty} d\lambda n\cdot A(\xi+\lambda n)\right] \ ,
\end{align}
along the light-cone direction $n^\mu$, as shown in Fig.~\ref{Fig:tmdpdf}.The $W_\perp$ is a transverse gauge-link at light-cone infinity to maintain gauge-invariance..
If one uses LFQ and ignores the transverse gauge-link, the above distribution
is just $\langle P|b^\dagger (x,\vec{k}_\perp)b(x,\vec{k}_\perp)|P\rangle$ for $x>0$, as expected.


\begin{figure}[htb]
	\centering
	\includegraphics[width=0.9\columnwidth]{images/tmd-TMDDY_fig10.pdf}
	\caption{The space-time picture of TMDPDF for DY and SIDIS process. The circled crosses denote the quark-link vertices. Notice that the vertices are placed at $\lambda n+\vec{b}_\perp$ and $0$ which gives the same result as the symmetric choice in Eq.~(\ref{eq:beam}). }
	\label{Fig:tmdpdf}
\end{figure}

However, there are a number of qualifications in the above definition.
First, the light-like gauge-links ${\cal W}_{n}$ are chosen to be past-pointing in accordance with the
DY kinematics, but for SIDIS  they should be chosen as future-pointing, as shown in Fig.~\ref{Fig:tmdpdf}.
For unpolarized TMDPDFs there is no distinction between the two choices,
but for spin-dependent TMDPDFs there are physical consequences associated with
the direction of gauge-links.

Second, there exists a new type of divergence associated with the infinitely-long light-like gauge-links. These divergences are due to radiation of gluons
collinear to the light-like gauge-link and cannot be regularized by the standard UV regulators. An example is
the following integral in dimensional regularization (DR)~\cite{Ebert:2019okf},
\begin{equation}\label{eq:RD}
	I = \int dk^+dk^- \frac{f(k^+k^-)}{(k^+k^-)^{1+\epsilon}} = \frac{1}{2}
	\int \frac{dy}{y} \int dm^2 \frac{f(m^2)}{m^{2+2\epsilon}}\,,
\end{equation}
where $m^2=k^+k^-$ and $y = k^+/k^-$ is the rapidity-related variable. The divergences
in $y$ arise from large and small $y$ where the integral is unregulated. The contribution
from $k^+=0$ is called the light-zero mode in LFQ, where it is also called light-cone divergence
which causes considerable problems.

To regulate the light-cone or rapidity divergences, a number of methods have been introduced in the literature (for a review see~\cite{Ebert:2019okf}).
They can be put into two classes: on-light-cone regulators and off-light-cone regulators.
In the former case, the gauge-links are kept along the light-cone direction $n^\mu$ after regularization. For example, the so-called {\it $\delta$ regulator}~\cite{Echevarria:2015usa,Echevarria:2015byo} regularizes the gauge-link as:
\begin{align}
	&W_n(\xi)\rightarrow W_{n}(\xi)|_{\delta^-}\nonumber \\
	&={\cal P}{\rm exp}\left[-ig\int_0^{-\infty} d\lambda A^+(\xi+\lambda n)e^{-\frac{\delta^-}{2p^+} |\lambda|}\right] \,,
\end{align} and similarly for the conjugate direction. The $\delta$ regulator breaks gauge-invariance, but preserves the boost invariance $\delta^{\pm}\rightarrow e^{\pm Y}\delta^{\pm}$ where $Y$ is the rapidity of the Lorentz boost. Other on-light-cone regulators include the exponential regulator~\cite{Li:2016axz}, $\eta$ regulator~\cite{Chiu:2012ir}, analytical regulator~\cite{Becher:2010tm}, etc. In the remainder of this section, we will use the $\delta$ regulator as a representative whenever we need an on-light-cone regulator.

The off-light-cone regulator was introduced in ~\cite{Collins:1981uk,Ji:2004wu,Ji:2004xq,Collins:2011zzd}, and also used in ~\cite{Ji:2004wu}. This type of regulator chooses off-light-cone directions to avoid the rapidity divergence. One can choose, for instance, to deform the gauge-links into the space-like region:
\begin{align}
	n\rightarrow n_Y=n-e^{-2Y}\frac{p}{(p^+)^2}\,.
\end{align}
Here $Y$ plays the role of a rapidity regulator, as when $Y \rightarrow \infty$, $n_Y \rightarrow n$. In certain cases one can also deform $n_Y$ into time-like region~\cite{Collins:2004nx}.

The on-light-cone regulators are consistent with the spirit of parton physics, and therefore are useful to define COM-momentum-independent parton densities. The off-light-cone regulators, on the other hand, follow a similar spirit as LaMET, and therefore can be exploited for practical lattice QCD calculations, as we shall see in the next subsection.

To avoid light-cone divergences, from now on we include the rapidity regulator in the definition of the
light-cone TMDPDFs. Using the same label $f$ for the TMDPDFs in both momentum and
coordinate spaces, we have
\begin{align}
	&f(\lambda,b_\perp,\mu,\delta^-/P^+)\\
	&= \langle P| \bar \psi(\lambda n/2+\vec{b}_\perp)\slashed{n}{\cal W}_{n}(\lambda n/2+\vec{b}_\perp)|_{\delta^-}\psi(-\lambda n/2) |P\rangle  \,,\nn
\end{align}
where $\mu$ is the $\MS$ scale for UV renormalization.
Due to rotational invariance, the bare TMDPDF defined above is a function of $b_\perp=|\vec{b}_\perp|$, so we have omitted the vector arrow for $\vec{b}_\perp$ in $f$ and will do so throughout the discussion for the soft functions, quasi-TMDPDFs, etc. The subscript $\delta^-$ denotes that the staple-shaped gauge-link ${\cal W}$ is regulated by the $\delta$ regulator in the light-cone minus direction. $f$ diverges logarithmically as $\delta^- \rightarrow 0$, and the finite part also depends on the rapidity
regulator. To define the physical TMDPDF, we need to remove all divergences and rapidity regularization scheme dependencies in $f$, in a way similar to removing UV divergences in physical quantities.
\begin{figure}[htb]
	\centering
	\includegraphics[width=0.9\columnwidth]{images/tmd-soft_fig11.pdf}
	\caption{The soft function $S(b_\perp,\mu,\delta^+,\delta^-)$ as space-time Wilson-loop arising in the
		factorization of DY and SIDIS process.}
	\label{Fig:soft1}
\end{figure}

The rapidity divergence for TMDPDFs can be removed by the
soft function, which also plays an important role in TMD factorization.
Intuitively, the soft function represents a cross section for fast-moving
charged particles emitting soft gluons into final states. It has rapidity divergence associated with the light-cone direction, which
is ultimately related to the mass singularity. The TMD soft function corresponding to Drell-Yan process is defined~\cite{Collins:2011ca,Echevarria:2015byo} as
\begin{align}\label{eq:soft}
	& S(b_\perp,\mu,\delta^+,\delta^-)\nonumber \\
	&=\frac{{\rm Tr}\langle 0|{\cal \bar T}W_p(\vec{b}_\perp)|_{\delta^+}W_n^{\dagger}(\vec{b}_\perp)|_{\delta^-}{\cal T} W_n(0)|_{\delta^-}W^{\dagger}_p(0)|_{\delta^+}|0\rangle }{N_c} \nonumber \\
	&=\frac{{\rm tr}\langle 0|{\cal W}_{n}(\vec{b}_\perp)|_{\delta^+} {\cal W}^{\dagger}_{p}(\vec{b}_\perp)|_{\delta^-}|0\rangle}{N_c} \,,
\end{align}
where ${\cal T}/\bar{\cal T}$ stands for time/anti-time ordering. The first equality defines the soft function in terms of cut-diagrams as an amplitude square. Since the soft function for DY process is independent of time ordering, one can also define it with a single time ordering or no time ordering, leading to the second equality. The staple-shaped gauge-link ${\cal W}_n$ is defined in Eq.~(\ref{eq:staplen}),
while the staple-shaped gauge-link ${\cal W}_p$ is defined similarly as:
\begin{align}
	&{\cal W}_p(\xi)=W^{\dagger}_{p}(\xi)W_{\perp}W_{p}(0) \label{eq:staplep} \ , \\
	&W_{p}(\xi)= {\cal P}{\rm exp}\left[-ig\int_{0}^{-\infty} d\lambda p\cdot A(\xi+p \lambda)\right] \ .
\end{align}
The soft function is shown in Fig.~\ref{Fig:soft1} as a Wilson loop in Minkowski space.

If the rapidity divergences are multiplicative, one can use $S$ as the rapidity renormalization factor for the TMDPDF defined in \eq{beam}. In on-light-cone schemes such as the $\delta$ regularization, it has been argued in~\cite{Vladimirov:2017ksc} based on conformal transformation that the rapidity divergences are indeed multiplicative in position space. For each of the staple-shaped light-like gauge-links, the rapidity divergence is proportional to $\exp\left[-(1/2)K(b_\perp,\mu)\ln \left(\mu^2/2(\delta^{\pm})^2\right)\right]$ where $K(b_\perp,\mu)$ is the nonperturbative Collins-Soper evolution kernel~\cite{Collins:1981uk}. Thus at small $\delta^{\pm}$, we can write
\begin{equation}
	S(b_\perp,\mu,\delta^+,\delta^-)=e^{\ln \frac{\mu^2}{2\delta^+\delta^-}K(b_\perp,\mu)+{\cal D}_2(b_\perp,\mu)} \ \ ,
\end{equation}
where ${\cal D}_2(b_\perp,\mu)$ is a $b_\perp$-dependent but rapidity-independent function. Notice that our definitions of $\delta^{\pm}$ differ from those in Ref.~\cite{Echevarria:2015byo} by a factor of $\sqrt{2}$ due to our normalization of light-cone vectors.

The soft-function in $\delta$ regularization satisfies the renormalization group equation
\begin{align}
	&\mu^2\frac{d}{d\mu^2}\ln S(b_\perp,\mu,\delta^+,\delta^-)\nonumber \\
	&=-\Gamma_{\rm cusp}(\alpha_s)\ln \frac{\mu^2}{2\delta^+\delta^-} + \gamma_s(\alpha_s)\ ,
\end{align}
where $\Gamma_{\rm cusp}(\alpha_s) $ is the light-like cusp anomalous dimension~\cite{Polyakov:1980ca,Korchemsky:1987wg} and the $\gamma_s(\alpha_s)$ is the soft anomalous dimension~\cite{Korchemskaya:1992je}. The Collins-Soper kernel and the rapidity-independent part ${\cal D}_2$ satisfy the RGEs:
\begin{align}
	\mu^2\frac{d}{d\mu^2} K(b_\perp,\mu)=-\Gamma_{\rm cusp}(\alpha_s) \label{eq:cusp}\ ,\\
	\mu^2\frac{d}{d\mu^2} {\cal D}_2(b_\perp,\mu)=\gamma_s(\alpha_s)-K(b_\perp,\mu) \label{eq:ad_d2}\ .
\end{align}

At one-loop, the soft function $S(b_\perp,\mu,\delta^+,\delta^-)$ is given by~\cite{Echevarria:2012js} :
\begin{align}
	&S(b_\perp,\mu,\delta^+,\delta^-)\nonumber \\ &=1+\frac{\alpha_s C_F}{2\pi}\left(L_b^2-2L_b \ln \frac{\mu^2}{2\delta^+\delta^-}+\frac{\pi^2}{6}\right) \ ,
\end{align}
where $L_b=\ln\left(\mu^2 b_\perp^2e^{2\gamma_E}/4\right)$.
Therefore, we have at the leading order,
\begin{align}
	K(b_\perp,\mu)&=-\frac{\alpha_sC_F}{\pi} L_b  \,,\\
	{\cal D}_2(b_\perp,\mu)&=\frac{\alpha_s C_F}{2\pi}\left(L_b^2+\frac{\pi^2}{6}\right) \ ,
\end{align}
and $\Gamma_{\rm cusp}=\alpha_s C_F/\pi +{\cal O}(\alpha_s^2)$, $\gamma_s={\cal O}(\alpha_s^2)$. It is worth pointing out that $K$~\cite{Li:2016ctv,Vladimirov:2016dll} and ${\cal D}_2$~\cite{Li:2016ctv} are known to 3-loop order in the exponential regularization scheme.

With the above soft function, we can take its square root to perform rapidity renormalization for the bare TMD correlator. The square root can be explained as follows: $S$ contains two staples, while $f$ contains one, thus the rapidity divergences as well as scheme dependencies in $S$ are twice as those in $f$. This leads to the following definition of the renormalized physical TMDPDF~\cite{Echevarria:2012js,Collins:2012uy}:
\begin{align}\label{eq:TMD}
	f^{\rm TMD}(x,b_\perp,\mu,\zeta)=\lim_{\delta^- \rightarrow 0}\frac{f(x,b_\perp,\mu,\delta^-/P^+)}{\sqrt{S(b_\perp,\mu,\delta^- e^{2y_n},\delta^-)}} \ ,
\end{align}
where the rapidity scale reads
\begin{align} \label{zeta}
	\zeta=2(xP^+)^2 e^{2y_n} \ .
\end{align}
The rapidity dependence in the numerator of the right-hand side of Eq.~(\ref{eq:TMD}) has the form $\exp[-\frac{1}{2}K(b_\perp,\mu)\ln \frac{(\delta^-)^2}{(xP^+)^2}]$, while in the denominator it behaves as $\exp[\frac{1}{2}K(b_\perp,\mu)\ln \frac{\mu^2}{2(\delta^-)^2e^{2y_n}}]$. The $\delta^-$ dependence thus cancels out in the ratio, leaving a dependence on the rapidity scale as $\exp[-\frac{1}{2}K(b_\perp,\mu)\ln \frac{\mu^2}{2(xP^+)^2e^{2y_n}}]$, which is controlled by the so-called Collins-Soper evolution equation:
\begin{align}\label{eq:cs}
	2\zeta \frac{d}{d\zeta} \ln  f^{\rm TMD}(x,b_\perp,\mu,\zeta)=K(b_\perp,\mu) \ .
\end{align}
The $\zeta$-dependence comes from the initial-state quark radiation and is intrinsically nonperturbative for large $b_\perp$. $f^{\rm TMD}(x,b_\perp,\mu,\zeta)$ is the standard object to be matched to in LaMET.

We should emphasize that although $f^{\rm TMD}$ is free from rapidity divergences, it does contain soft radiation from the
charged particles in the initial state. This can be seen clearly by considering Feynman diagrams for the unsubtracted $f$
and applying soft approximation to gluons. ``One-half'' of the soft contribution in $f$
is subtracted to define the physical $f^{\rm TMD}$ due to the requirement of factorization of
physical processes. The remaining soft radiation also has a natural rapidity cut-off
associated with $\ln(xP^+)$, reflected in the $\zeta$-dependence. What is remarkable,
however, is that $f^{\rm TMD}$ is rapidity-regulator independent. Although a general proof to all
orders in perturbation theory is beyond the scope of this review, it is due to
factorization and exponentiation of the soft physics in $f$ and thus the scheme cancellation
can be done systematically in the exponent. It worth mentioning that in old-fashioned or SCET-like approaches, one can define the ``subtracted'' TMDPDF or ``beams functions'' that contains only collinear physics. However, they are generically scheme dependent and must be combined with an extra soft functions in factorization theorems. At one-loop level, the scheme-independent one-loop TMDPDF for an external quark state reads,
\begin{align}\label{eq:ftmd_1loop}
	&f^{\rm TMD}(x,b_\perp,\mu,\zeta) = \delta(1-x)\nonumber \\
	&+\frac{\alpha_sC_F}{2\pi}F(x,\epsilon_{\tiny\mbox{IR}},b_\perp,\mu)\theta(x)\theta(1-x)+\frac{\alpha_sC_F}{2\pi}\delta(1-x)\nonumber \\
	&\times \left[-\frac{1}{2}L_b^2+\left(\frac{3}{2}-\ln \frac{\zeta}{\mu^2}\right)L_b+\frac{1}{2}-\frac{\pi^2}{12}\right] \ ,
\end{align}
where
\begin{align}
	F(x,\epsilon_{\tiny\mbox{IR}},b_\perp,\mu)=\left[-\left(\frac{1}{\epsilon_{\tiny\mbox{IR}}}+L_b\right)\frac{1+x^2}{1-x}+1-x\right]_+\,.
\end{align}
Two-loop order results for the TMDPDFs can be found in~\cite{Catani:2011kr,Catani:2012qa,Gehrmann:2014yya,Luebbert:2016itl,Echevarria:2016scs,Luo:2019hmp} and three-loop order results can be found in~\cite{Luo:2019szz}.

The physical TMDPDF also satisfies the RG equation,
\begin{align}
	\gamma_\mu(\mu,\zeta)&=\mu^2\frac{d}{d\mu^2}\ln f^{\rm TMD}(x,b_\perp,\mu,\zeta) \nonumber \\
	&\equiv \frac{1}{2}\Gamma_{\rm cusp}(\alpha_s)\ln \frac{\mu^2}{\zeta}-\gamma_H(\alpha_s)  \ ,
\end{align}
where $\gamma_H$ is called the hard anomalous dimension. At one-loop, the cusp and hard anomalous dimensions read
\begin{align}\label{eq:ad_hard}
	\Gamma_{\rm cusp}(\alpha_s) =\frac{\alpha_s C_F}{\pi}; ~~~~~
	\gamma_H(\alpha_s) =-\frac{3\alpha_sC_F}{4\pi}  \ .
\end{align}
Recently the cusp anomalous dimension have been calculated to 4-loops~\cite{Henn:2019swt,vonManteuffel:2020vjv}.

Combining the RGE and the rapidity evolution equation for the TMDPDF, one obtains the consistency condition :
\begin{align}
	\mu^2 \frac{d}{d\mu^2} K(b_\perp,\mu)
	&= -2 \zeta \frac{d}{d\zeta} \gamma_\mu(\mu,\zeta) =- \Gamma_{\rm cusp}(\alpha_s(\mu)) \ ,
\end{align}
from which one finds a resummed form for the Collins-Soper kernel:
\begin{align}
	K(b_\perp,\mu) &= -2 \int_{1/b_\perp}^\mu \frac{d\mu'}{\mu'} \Gamma_{\rm cusp}(\alpha_s(\mu')) + K(\alpha_s(1/b_\perp)) \ .
\end{align}
Here $K(\alpha_s(1/b_\perp))$ contains both perturbative and non-perturbative contributions. The TMDPDFs at different scales are then related by
\begin{align} \label{eq:tmd_evolution}
	&f^{\rm TMD}(x, b_\perp, \mu, \zeta) = f^{\rm TMD}(x, b_\perp, \mu_0, \zeta_0) \\
	& \times\exp\biggl[ \int_{\mu_0}^\mu \frac{d\mu'}{\mu'} \gamma_\mu(\mu',\zeta_0) \biggr]
	\exp\biggl[ \frac12 K(b_\perp,\mu) \ln\frac{\zeta}{\zeta_0} \biggr]  \ . \nn
\end{align}
The double-scale evolution in the $\mu-\zeta$ plane for phenomenology has been recently studied in~\cite{Scimemi:2018xaf}.
With the scheme-independent physical TMDPDF defined above, the DY cross section at $s=(P_A+P_B)^2$ and small $Q_\perp$ can be factorized as
\begin{align}
	&\frac{d\sigma}{dQ_\perp^2}=\int dx_Adx_B d^2b_\perp e^{i\vec{b}_\perp \cdot \vec{Q}_\perp}\hat \sigma(x_Ax_Bs,\mu)\nonumber \\ &\times f_A^{\rm TMD}(x_A,b_\perp,\mu,\zeta_A)f_B^{\rm TMD}(x_B,b_\perp,\mu,\zeta_B)+ ... \ .
\end{align}
The rapidity scales satisfy $\zeta_A\zeta_B= Q^4\equiv(x_Ax_Bs)^2$. The remaining term at large but finite $Q^2$ are called power corrections or ``higher-twist'' contributions. A detailed study of the power corrections to TMD factorization is beyond the scope of this review. Without mention we will omit all the power-corrections in equations. The QCD part of the hard cross section $\hat{\sigma}$ at one-loop level reads
\begin{align}
	\hat \sigma(x_A,x_B)=\left|1+\frac{\alpha_sC_F}{4\pi}\left(-L^2_{Q}+3L_Q-8+\frac{\pi^2}{6}\right)\right|^2 \ ,
\end{align}
where $L_Q=\ln \frac{-Q^2-i0}{\mu^2}$, and the result is now known up to three loops (see~\cite{Moch:2005tm,Baikov:2009bg,Lee:2010cga,Gehrmann:2010ue} and the references therein). Similarly for the SIDIS process we have
\begin{align}
	\frac{d\sigma}{dQ_\perp^2}&=\int dxdz d^2b_\perp e^{i\vec{b}_\perp \cdot \vec{Q}_\perp}H(x,z,\mu,Q)\nonumber \\
	&\quad \times f^{\rm TMD}(x,b_\perp,\mu,\zeta_A)d^{\rm TMD}(z,b_\perp,\mu,\zeta_B) \ ,
\end{align}
where $d^{\rm TMD}(z,b_\perp,\mu,\zeta_B)$ is the TMD fragmentation function and $H$ is the hard kernel.

%------------------------------------------------------------------------------


\subsection{Lattice Quasi-TMDPDFs and Matching}
\label{sec:quasitmd}
%------------------------------------------------------------------------------S
Before LaMET, there had been efforts to access TMD physics from lattice QCD by calculating the ratios of the $x$-moments of TMDPDFs~\cite{Hagler:2009mb,Musch:2010ka,Musch:2011er,Engelhardt:2015xja,Yoon:2017qzo}, which are free from complications associated with the soft function and can be compared to certain experimental observables. In LaMET, we are more interested in obtaining the full $x$ and $\vec{k}_\perp$ dependence of the TMDPDFs~\cite{Ji:2014hxa,Ji:2018hvs,Ebert:2018gzl,Ebert:2019okf,Ji:2019sxk,Ji:2019ewn}. Therefore, a proper treatment of the soft function subtraction and matching is essential.
The earliest suggestion of a bent soft function in~\cite{Ji:2014hxa} and the follow-up work~\cite{Ebert:2019okf} has the correct IR logarithms at one-loop order, but this is expected to break down at higher-loop orders~\cite{Ji:2019sxk}, thus not allowing for a perturbative matching. Another suggestion which uses a naive rectangle-shaped Wilson loop~\cite{Ji:2018hvs,Ebert:2019okf} does not possess the correct IR physics, either. Nevertheless, in~\cite{Ebert:2018gzl} important progress was made for calculating the nonperturbative
Collins-Soper kernel $K(b_\perp,\mu)$ from the ratio of quasi-TMDPDFs at two different large momenta.
Recently, some of the authors showed~\cite{Ji:2019sxk,Ji:2019ewn} that the quasi-TMDPDF combined with a reduced soft function capture the correct IR physics to all-orders and thus allow for a perturbative matching to the physical TMDPDF.

To construct such quasi-TMDPDFs, the collinear part can be treated in a way similar to the collinear PDFs, while the soft
piece is more challenging. Our starting point is that the physical $f^{\rm TMD}$ is independent of the rapidity regulator, so one can use a scheme in which the gauge-links in both $f$ and $S$
are off the light-cone, such as that used in~\cite{Collins:2011zzd}. In this case, one can use Lorentz symmetry to
boost the staple-shaped gauge-link ${\cal W}_n$ in $f$ to a purely space-like staple
with no time dependence. However, one can only use
this trick for one of the staples in $S$, say ${\cal W}_n$, whereas the other one
${\cal W}_p$ is still time-dependent. In other words, there is no way to get rid of the time dependence
in $S$ entirely with Lorentz boost alone. This is natural because $S$ in fact represents the square of
an $S$-matrix, which appears to be intrinsically Minkowskian. However, using the LaMET principle that time
dependence of an operator can be simulated through external physical states at large momentum, we find that
$S$ can indeed be calculated on the lattice in the off-light-cone scheme as a form factor. A detailed discussion will be given in
the next subsection. Here we assume that this is true, and discuss the matching between
quasi- and physical TMDPDFs.

First, we define the quasi-TMDPDF with a staple–shaped gauge-link along the $z$ direction~\cite{Ji:2014hxa,Ji:2018hvs,Ebert:2019okf,Ji:2019ewn} as
\begin{align}\label{eq:quasi_TMD}
	& \tilde f(\lambda ,b_\perp,\mu,\zeta_z) \\
	&=\! \lim_{L \rightarrow \infty}  \frac{\langle P| \bar \psi\big(\frac{\lambda n_z }{2}\!+\!\vec{b}_\perp\big)\gamma^z{\cal W}_{z}(\frac{\lambda n_z}{2}\!+\!\vec{b}_\perp;L)\psi\big(\!-\!\frac{\lambda n_z}{2}\big) |P\rangle}{\sqrt{Z_E(2L,b_\perp,\mu)}} \ , \nonumber
\end{align}
where the $\MS$ renormalization is implied, and
\begin{align}\label{eq:staplez}
	&{\cal W}_z(\xi;L)=W^{\dagger}_{z}(\xi; L)W_{\perp}W_{z}(-\xi^zn_z;L)  \ ,\\
	&W_{z}(\xi;L)= {\cal P}{\rm exp}\Big[-ig\int_{\xi^z}^{L} d\lambda n_z\cdot A(\vec{\xi}_\perp\!+\!n_z \lambda)\Big] \ .
\end{align}
Here $\xi^z=-\xi\cdot n_z$ and $\zeta_z=(2xP^z)^2$ is the Collins-Soper scale of the quasi-TMDPDF. $W_{\perp}$ is inserted at $z=L$ to maintain explicit gauge invariance. $\sqrt{Z_E(2L,b_\perp,\mu,0)}$ is the square root of the vacuum expectation value of a flat rectangular Euclidean Wilson-loop along the $n_z$ direction with length $2L$ and width $b_\perp$:
\begin{align}\label{eq:Z_E}
	Z_E(2L,b_\perp,\mu)=\frac{1}{N_c}{\rm Tr}\langle 0|W_{\perp}{\cal W}_z(\vec{b}_\perp;2L)|0\rangle \ .
\end{align}
Again, $\gamma^z$ can be replaced by $\gamma^t$ as in the collinear quasi-PDF.
For a depiction of $\tilde f$ and $Z_E$ see Fig.~\ref{fig:quasiTMD}.

The purpose of the factor $Z_E$ is as follows. At large $L$, the naive quasi-TMD correlator in the numerator of Eq.~(\ref{eq:quasi_TMD}) contains divergences that go as $e^{-LE(b_\perp,\mu)}$ where $E(b_\perp)$ is the ground state energy of a pair of static heavy-quarks. $E(b_\perp,\mu)=2\delta m+ V(b_\perp,\mu)$ contains both the linear divergent mass corrections $2\delta m$ and the heavy-quark potential $V(b_\perp,\mu)$ due to mutual interactions. In literature the $LV(b_\perp,\mu)$ part was sometimes called the ``pinch pole singularity.'' Therefore, we introduce the square root of a rectangular Wilson-loop $Z_E(2L,b_\perp,\mu)$ with twice the length to cancel all these divergences and guarantee the existence of the $L\to\infty$ limit after the subtraction. The introduction of $\sqrt{Z_E}$ also removes additional contributions from the transverse gauge link. An alternative approach to avoid the pinch-pole singularity was proposed in~\cite{Li:2016amo}. We should mention that although the $\sqrt{Z_E}$ subtraction removes all the linear divergences, the logarithmic UV divergences are still present. Therefore, a non-perturbative renormalization of $\tilde f$ on the lattice is still required, which has been studied in the RI/MOM scheme~\cite{Shanahan:2019zcq}, and its matching to the $\MS$ scheme has been calculated at one-loop order~\cite{Constantinou:2019vyb,Ebert:2019tvc}.
\begin{figure}[htb]
	\centering
	\includegraphics[width=0.8\columnwidth]{images/tmd-quasiTMD_fig12.pdf}.
	\caption{The quasi-TMDPDF (upper) and the Euclidean Wilson-loop $Z_E(2L,b_\perp,\mu,0)$ (lower). In the figure, $A=\lambda n_z/2+\vec{b}_\perp/2$, $B=-\lambda n_z/2-\vec{b}_\perp/2$ and $C=Ln_z+{\vec b}_\perp$. The crosses denote the quark-link vertices.}
	\label{fig:quasiTMD}
\end{figure}


The quasi-TMDPDFs defined above satisfy the following RGE~\cite{Collins:1981uk,Ji:2014hxa,Ji:2019ewn}
\begin{align}
	\mu^2 \frac{d}{d\mu^2}\ln \tilde f(x,b_\perp,\mu,\zeta_z)=\gamma_{F}(\alpha_s(\mu)) \ ,
\end{align}
where $\gamma_F$ is the anomalous dimension for the heavy-to-light current in \sec{renorm-nl}. This is due to the fact that the quasi-TMDPDF, after the self-energy subtraction, contains only logarithmic UV divergences associated with quark-Wilson-line vertices. In the $\MS$ scheme, the one-loop quasi-TMDPDF in an external quark state with momentum $(p^z,0,0,p^z)$ reads~\cite{Ji:2018hvs,Ebert:2019okf}
\begin{align}
	&\tilde f(x,b_\perp,\mu,\zeta_z)=\nonumber \\ &1+\frac{\alpha_sC_F}{2\pi}F(x,\epsilon_{\tiny \mbox{IR}},b_\perp,\mu)\theta(x)\theta(1-x)+\frac{\alpha_sC_F}{2\pi}\delta(1-x) \nonumber \\
	& \times \Bigg[-\frac{1}{2}L_b^2+L_b\Big(\frac{5}{2}-L_z\Big)-\frac{3}{2}-\frac{1}{2}L_z^2+L_z\Bigg]\ ,
\end{align}
where $L_z=\ln(\zeta_z/\mu^2)$. As expected, the $L$ dependence has been cancelled in the large $L$ limit.

As there is no light-like gauge-link in $\tilde f$, no additional rapidity regulator is needed. Instead, there is an explicit dependence on the hadron momentum (or energy), which is similar to the momentum RGE for collinear quasi-PDF. The momentum (rapidity) evolution equation for $\tilde f$ reads~\cite{Collins:1981uk,Ji:2014hxa,Ji:2019ewn},
\begin{align}
	P^z\frac{d}{dP^z}\ln \tilde f(x,b_\perp,\mu,\zeta_z) \!=\! K(b_\perp,\mu)\!+\!{\cal G}\Big(\frac{(P^z)^2}{\mu^2}\Big) \ ,
\end{align}
where ${\cal G}(\zeta_z/\mu^2)$ is perturbative and $K(b_\perp,\mu)$ is the Collins-Soper kernel. A similar equation was proven for off-light-cone TMD-fragmentation functions in~\cite{Collins:1981uk}. From this equation, it is clear that a correct matching to $f^{\rm TMD}(x,b_\perp,\mu,\zeta)$ with arbitrary $\zeta$ must include $K(b_\perp,\mu)$ to compensate the $P^z$ dependence.

There is actually one more requirement for the matching: there is a rapidity scheme dependence which must
be removed, since the quasi-TMDPDF can be viewed as defined with an off-light-cone regulator along the $z$ direction. To understand this dependence, let us consider $f$ again in the off-light-cone regularization, where there are rapidity divergences. The divergence is cancelled by the square root of an off-light-cone soft function $S_{\rm DY}(b_\perp,\mu,Y,Y')$, with $Y,Y'$ being the rapidities of the off-light-cone space-like vectors $p\rightarrow p_Y= p-e^{-2Y}(p^+)^2n $ and $n \rightarrow n_{Y'}=n-e^{-2Y'}p/(p^+)^2$. Schematically, we have:
\begin{align}\label{eq:S_DY}
	S_{\rm DY}(b_\perp,\mu,Y,Y')= \frac{{\rm tr}\langle 0|{\cal W}_{n_{Y'}}(\vec{b}_\perp) {\cal W}^{\dagger}_{p_Y}(\vec{b}_\perp)|0\rangle}{N_c\sqrt{Z_E}\sqrt{Z_E}}  \ ,
\end{align}
where ${\cal W}_{n_{Y'}}(\vec{b}_\perp)$ and ${\cal W}^{\dagger}_{p_Y}(\vec{b}_\perp)$ are staple-shaped gauge-links in $n_{Y'}$, $p_Y$ directions, respectively. $\sqrt{Z_E}$ is introduced to subtract the pinch pole singularities for the off-light-cone staple-shaped gauge-links. In terms of $\ln \rho^2=2(Y+Y')$ sometimes we also write this soft function as $S_{\rm DY}(b_\perp,\mu,\rho)$. At large $\rho $, we have
\begin{align}
	S_{\rm DY}(b_\perp,\mu,Y,Y')= e^{(Y+Y') K(b_\perp,\mu)+{\cal D}(b_\perp,\mu)}+... \,.
\end{align}
We can perform a Lorentz boost of ${\cal W}_{n_{Y'}}(\vec{b}_\perp) {\cal W}^{\dagger}_{p_Y}(\vec{b}_\perp)$ in \eq{S_DY} such that one of the gauge-links, say ${\cal W}_{n_{Y'}}$, is boosted to the equal-time version ${\cal W}_z$ in $\tilde f$, whereas the other gauge-link ${\cal W}_{n_{Y}}$ is boosted to ${\cal W}_{n_{Y+Y'}}$. The soft function becomes $S_{\rm DY}(b_\perp,\mu,Y+Y',0)$ which
contains light-cone divergence for the $p_{Y+Y'}$ direction, but is still the same $S_{\rm DY}(b_\perp,\mu,Y,Y')$ due to boost invariance. The square-root of the finite part $e^{{\cal D}(b_\perp,\mu)}$
is exactly what is needed to cancel the rapidity-scheme dependence. We define the rapidity-independent part as
the reduced soft function:
\begin{align}\label{eq:reduced}
	S_r(b_\perp,\mu)\equiv e^{-{\cal D}(b_\perp,\mu)} \,.
\end{align}
Based on the renormalization property of non-light-like Wilson-loops, the reduced soft function satisfies the RG equation
\begin{align}
	\mu^2 \frac{d}{d\mu^2} \ln S_r(b_\perp,\mu)=\Gamma_S(\alpha_s) \,,
\end{align}
where $\Gamma_S$ is the constant part of the cusp-anomalous dimension at large hyperbolic cusp angle $Y+Y'$ for the off-light-cone soft function:
\begin{align}
	&\mu^2 \frac{d}{d\mu^2} \ln S_{\rm DY}(b_\perp,\mu,Y,Y')\nonumber \\
	&\qquad\qquad=-(Y+Y')\Gamma_{\rm cusp}(\alpha_s)-\Gamma_S(\alpha_s)\,.
\end{align}
At one-loop level~\cite{Ebert:2019okf},
\begin{align}
	S^{(1)}_{\rm DY}(b_\perp,\mu,Y,Y')=\frac{\alpha_sC_F}{2\pi}\big[2-2(Y+Y')\big]L_b \ ,
\end{align}
and $\Gamma_S^{(1)} (\alpha_s)= \alpha_s C_F/\pi$.
Based on RGE, at two-loop level ${\cal D}(b_\perp,\mu)$ can be predicted to be
\begin{align}
	{\cal D}^{(2)}(b_\perp,\mu)=c_2+\Gamma_S^{(2)}L_b-\frac{\alpha_s^2\beta_0C_F}{2\pi}L_b^2 \ ,
\end{align}
where
\begin{align}
	\Gamma_S^{(2)}=-\frac{\alpha_s^2}{\pi^2}\Big[C_FC_A\big(-\frac{49}{36}+\frac{\pi^2}{12}-\frac{\zeta_3}{2}\big)+C_FN_F\frac{5}{18}\Big]\nn
\end{align}
is the two-loop anomalous dimension for $S_r$ which can be extracted from~\cite{Grozin:2015kna}, $\beta_0=-\left(\frac{11}{3}C_A-\frac{4}{3}N_f T_F\right)/(2\pi)$ is the coefficient of one-loop $\beta$-function, and $c_2$ is a constant to be determined by explicit calculation.

After taking into account the reduced soft function, we can now write down
the matching formula between the quasi-TMDPDF and the scheme-independent TMDPDF\cite{Ji:2019ewn}:
\begin{align}\label{eq:quasiTMDmatch}
	&f^{\rm TMD}(x,b_\perp,\mu,\zeta)\\
	&=H\left(\frac{\zeta_z}{\mu^2}\right) e^{-\ln (\frac{\zeta_z}{\zeta})K(b_\perp,\mu)}\tilde f(x,b_\perp,\mu,\zeta_z)S_r^{\frac{1}{2}}(b_\perp,\mu)+...\ ,\nn
\end{align}
where the power-corrections of order ${\cal O}\left(\Lambda_{\rm QCD}^2/\zeta_z,M^2/(P^z)^2,1/(b^2_\perp\zeta_z) \right)$.
The above relation except for the definition of $S_r(b_\perp,\mu)$ was argued to hold in~\cite{Ebert:2019okf}, where the unknown function $g_q^S$ in Eq.~(5.3) should be identified as the reduced soft function here; it has also been confirmed recently in~\cite{Vladimirov:2020ofp}.


We now explain the individual factors of the formula.

\begin{enumerate}
	\item The factor $ H(\zeta_z/\mu^2)$ is the perturbative matching kernel, which is a function of $\zeta_z/\mu^2=(2xP^z)^2/\mu^2$. The kernel is responsible for the large logarithms of $P^z$ generated by the ${\cal G}(\zeta_z/\mu^2)$ term of the momentum RG equation. Unlike the case of quasi-PDFs, the momentum fractions of the quasi-TMDPDF and the TMDPDF are the same. This is due to the fact that at leading power in $1/\zeta_z$ expansion, the $k_\perp$ integral is naturally cut off by the transverse separation around $k_\perp \sim 1/b_\perp \ll P^z$. Therefore, the momentum fraction can only be modified by collinear modes for which there are no distinction between $x=k^z/P^z$ and $x=k^+/P^+$. In comparison, for the $\vec{k}_\perp$ integrated quasi-PDF, the $k_\perp \ge P^z$ region leads to non-trivial $x$ dependence outside the physical region. This is also consistent with the fact that the momentum evolution equation for quasi-TMDPDF is local in $x$ instead of being a convolution.
	\item The factor $\exp\big[\ln (\frac{\zeta_z}{\zeta})K(b_\perp,\mu)\big]$ is the part involving the Collins-Soper evolution kernel. From the momentum evolution equation, it is clear that at large $P^z$ there are logarithms of the form $K(b_\perp,\mu)\ln \frac{\zeta_z}{\mu^2}  $ with $\zeta_z$ being the natural Collins-Soper scale. Therefore, to match to the TMDPDF at arbitrary $\zeta$, a factor $\exp\big[\ln (\frac{\zeta_z}{\zeta})K(b_\perp,\mu)\big]$ is required to compensate the difference. An important implication of this property is that one can obtain the Collins-Soper kernel $K(b_\perp,\mu)$ by constructing the ratio of quasi-TMDPDFs at two different momenta or $\zeta_z$'s~\cite{Ebert:2018gzl},
	\begin{align}
		\label{eq:TMD-CS_kernel}
		\ \ \quad\frac{\tilde f(x,b_\perp,\mu,\zeta_{z,1})}{\tilde f(x,b_\perp,\mu,\zeta_{z,2})}=\frac{H\left(\frac{\zeta_{z,2}}{\mu^2}\right)}{H\left(\frac{\zeta_{z,1}}{\mu^2}\right)}\left(\frac{\zeta_{z,1}}{\zeta_{z,2}}\right)^{K(b_\perp,\mu)} \ .
	\end{align}
	Thus given the $\tilde f$'s at the two rapidity scales, the Collins-Soper kernel $K(b_\perp)$ can be obtained.
	
\end{enumerate}

Combining the RGEs of the quasi-TMDPDF $\tilde f$, reduced soft function $S_r$ and physical TMDPDF $f^{\rm TMD}$, we obtain the RGE of the matching kernel $H\left(\frac{\zeta_z}{\mu^2}\right)$ \cite{Ji:2019ewn},
\begin{align}
	\mu^2 \frac{d}{d\mu^2} \ln H^{-1}\left(\frac{\zeta_z}{\mu^2}\right)=\frac{1}{2}\Gamma_{\rm cusp}(\alpha_s)  \ln\frac{\zeta_z}{\mu^2}\!+\!\frac{\gamma_C(\alpha_s)}{2}  \ ,
\end{align}
where $\gamma_C(\alpha_s)=2\gamma_F(\alpha_s)+\Gamma_S(\alpha_s)+2\gamma_H(\alpha_s)$. The matching kernel is closely related to the perturbative part of the rapidity evolution kernel ${\cal G}\left(\frac{\zeta_z}{\mu^2}\right)$ through
\begin{align}
	2\zeta_z\frac{d}{d\zeta_z}\ln H^{-1}\left(\frac{\zeta_z}{\mu^2}\right)={\cal G}\left(\frac{\zeta_z}{\mu^2}\right) \ .
\end{align}
Again, we can see that the anomalous dimension of ${\cal G}\left(\frac{\zeta_z}{\mu^2}\right)$ is $\Gamma_{\rm cusp}(\alpha_s)$.

It is convenient to write $H$ in the exponential form, $H=e^{-h}$. Collecting all the above results, one obtains at one-loop level~\cite{Ji:2018hvs,Ebert:2019okf}
\begin{align}
	h^{(1)}\left(\frac{\zeta_z}{\mu^2}
	\right)=\frac{\alpha_sC_F}{2\pi}\left(-2+\frac{\pi^2}{12}-\frac{L_z^2}{2}+L_z\right) \ .
\end{align}
Similar as before, the two loop contribution $h^{(2)}$ is predicted to be
\begin{align}
	&h^{(2)}\left(\frac{\zeta_z}{\mu^2}\right)=c'_2-\frac{1}{2}\left(\gamma^{(2)}_C-\alpha_s^2\beta_0 c_1\right)\ln\frac{\zeta_z}{\mu^2}\\
	&-\frac{1}{4}\left(\Gamma^{(2)}_{\rm cusp}-\frac{\alpha_s^2\beta_0 C_F}{2\pi}\right)\ln^2\frac{\zeta_z}{\mu^2}-\frac{\alpha_s^2\beta_0C_F}{24\pi}\ln^3\frac{\zeta_z}{\mu^2} \ ,\nn
\end{align}
where $c_1=\frac{C_F}{2\pi}\left(-2+\frac{\pi^2}{12}\right)$ and $c'_2$ is again a constant to be determined in perturbation theory at two-loop level.

Finally, we compare the current formulation with previous approaches to lattice TMDPDF. First, we comment on the developments in~\cite{Hagler:2009mb,Musch:2010ka,Musch:2011er,Engelhardt:2015xja,Yoon:2017qzo} in which the $x$-moments of TMDPDF are extracted from ratio of quasi-TMDPDF.  From Eq.~(\ref{eq:quasiTMDmatch}), it is clear that both the matching kernel $H$ and the exponential factor of Collins-Soper kernel depends on $x$ non-trivially. Therefore, simply taking the ratio of moments for quasi-TMDPDF will not be sufficient to reproduce the same ratio for TMDPDF, although the soft function does cancel. This observation is also made recently in Ref.~\cite{Ebert:2020gxr}. Second, the quasi-TMDPDF defined with the naive rectangle-shaped soft function, i.e. $Z_E$, is $\tilde f$ in \eq{quasi_TMD}, so it is obvious that it still needs the reduced soft function $S_r$ to be matched to $f^{\rm TMD}$. As for the other proposal in~\cite{Ji:2014hxa,Ebert:2019okf}, it replaces $Z_E$ in $\tilde f$ with $S_{\rm bent}$ which is the vacuum matrix element of a spacelike bent-shaped Wilson loop with angle $\pi/2$ at each junction, and does not include the function $S^{1\over2}_r$ in \eq{quasiTMDmatch}. Although $\sqrt{S_{\rm bent}/Z_E}$ agrees with $S_r^{-1\over2}$ at one-loop order~\cite{Ji:2018hvs,Ebert:2019okf}, it is expected to be different at higher orders. In fact, for the anomalous dimension $\Gamma_{\pi\over2}$ defined through
\begin{align}
	\Gamma_{\frac{\pi}{2}}(\alpha_s) \equiv \mu^2\frac{d}{d\mu^2} \ln \left(\frac{S_{\rm bent}(L,b_\perp,\mu)}{Z_E(2L,b_\perp,\mu)}\right) \,,
\end{align}
it starts to deviate from $\Gamma_S(\alpha_s)$ at two-loop order~\cite{Grozin:2015kna}, as
\begin{align}
	-\Gamma_S(\alpha_s)&=\frac{\alpha_sC_F}{\pi} \\
	&+\frac{\alpha_s^2}{\pi^2}\left[C_FC_A\left(-\frac{49}{36}+\frac{\pi^2}{12}-\frac{\zeta_3}{2}\right)+C_FN_F\frac{5}{18}\right]  \,, \nn \\
	\Gamma_{\frac{\pi}{2}}(\alpha_s)&=\frac{\alpha_sC_F}{\pi} \\
	&+\frac{\alpha_s^2}{\pi^2}\left[C_FC_A\left(-\frac{49}{36}+\frac{\pi^2}{24}\right)+C_FN_F\frac{5}{18}\right] \,. \nn
\end{align}
In the equation, $\zeta_3=\sum_{n=1}^{\infty}(1/n^3) \ne \pi^2/12$, therefore the two anomalous dimensions are different. The differences in the anomalous dimension will result in different logarithmic behaviors in $b_\perp$, as the soft functions are dimensionless and depend on $b_\perp$ and $\mu$ only. At large $b_\perp$, it will lead to different IR physics that cannot be controlled by perturbation theory.

Combining the reduced soft function and the quasi-TMDPDF, one can effectively factorize the DY cross section,
\begin{align}
	&\sigma = \int dx_A dx_B d^2b_\perp e^{i\vec{Q}_\perp \cdot \vec{b}_\perp}\hat \sigma(x_A,x_B,Q^2,\mu)\nonumber \\ &\times \tilde f(x_A,b_\perp,\mu,\zeta_A) \tilde f(x_B,b_\perp,\mu,\zeta_B) S_r(b_\perp,\mu) \ .
\end{align}
where all non-perturbative quantities do not involve the light cone,
and can be calculated on lattice.

Spin-dependent TMDPDFs are also physically important. They can be computed
using LaMET theory~\cite{Ebert:2020gxr}. Again one can define quasi distributions just like the spin-independent ones.
For a general proton target $|PS\rangle$ and the general spin structure $\Gamma$ of the parton, the parent TMDPDF can be defined as :
\begin{align}\label{eq:parent_TMD}
	&f_{[\Gamma]}^{\rm TMD}(x,\vec{k}_\perp,\mu,\zeta)=\frac{1}{2P^+}\int\frac{d\lambda}{2\pi}\int\frac{d^2 \vec{b}_\perp}{(2\pi)^2}e^{-i\lambda x + i\vec{k}_\perp\cdot \vec{b}_\perp}\nonumber\\
	&\times \lim_{\delta^- \rightarrow 0}\frac{\langle PS|\bar\psi(\lambda n+\vec{b}_\perp)\Gamma\, {\cal W}_n(\lambda n +\vec{b}_\perp)|_{\delta^-}\psi(0)|PS\rangle}{\sqrt{S(b_\perp,\mu,\delta^-e^{2y_n},\delta^-)}} \,,
\end{align}
where the $\zeta=2(P^+)^2e^{2y_n}$ is the rapidity scale, see \sec{tmd} for more detail of the soft function subtraction. The individual spin-dependent TMD distributions can then be obtained through Lorentz decompositions ~\cite{Tangerman:1994eh,Ralston:1979ys,Mulders:1995dh}:
\begin{align}\label{eq:TMDPDF}
	f_{[\gamma^+]}^{\rm TMD}&=f_1-\frac{\epsilon^{ij}k^i S^j_\perp}{M}f_{1T}^\perp   \,,\\
	f_{[\gamma^+\gamma_5]}^{\rm TMD}&=S^+ g_1+\frac{\vec{k}_\perp\cdot \vec{S}_\perp}{M}g_{1T} \,, \\
	f_{[i\sigma^{i+}\gamma_5]}^{\rm TMD}&=S_\perp^i h_1+\frac{(2k^i k^j-\vec{k}_\perp^2\delta^{ij})S^j_\perp}{2M^2}h_{1T}^\perp\nn\\
	&\qquad+\frac{S^+ k^i}{M}h_{1L}^\perp+\frac{\epsilon^{ij} k^j}{M}h_1^\perp \,,
\end{align}
where we suppress the arguments $(x,\vec{k}_\perp,\mu,\zeta)$ in all distributions;
$f_1$, $g_1$, and $h_1$ are unpolarized, helicity and transversity TMDPDFs, respectively;
the indices $i$ and $j$ are in transverse space of $\vec{k}_\perp$; $S^+$ and $S^i_\perp$ are longitudinal and transverse spin components.

Note that the Sivers function $f^\perp_{1T}$~\cite{Sivers:1989cc} and the Boer-Mulders function $h^\perp_1$~\cite{Boer:1997nt} are $T$-odd. The orientation of the gauge-link have important effects on these two functions~\cite{Collins:2002kn,Collins:2011zzd}, such that they change sign between the DY and SIDIS processes.  In the light-cone gauge, these contributions arise from the transversal gauge-link at infinities~\cite{Belitsky:2002sm}. They are related to the phenomenologically interesting single transverse-spin asymmetry~\cite{Boer:1997nt,Efremov:2004tp,Collins:2005ie,Collins:2005wb}.

\subsection{Off-light-cone Soft Function}
\label{sec:offlight}
In previous subsections, the soft function has been introduced to define rapidity-scheme-independent
TMDPDFs. The major motivation of introducing the soft function is to capture nonperturbative effects due to
soft-gluon radiations from fast moving color-charges.
For many inclusive processes the soft radiations cancel in the total cross section, but for certain processes where a small transverse momentum is measured, such cancellation can be incomplete and result in measurable consequences.
In such cases, the TMD soft function is introduced to account for the soft-gluon effects
and appears in factorization theorems for the Drell-Yan (DY) process~\cite{Collins:1984kg,Collins:1988ig}
and semi-inclusive DIS (SIDIS)~\cite{Ji:2004wu,Ji:2004xq}.

To calculate the TMD physics nonperturbatively, formulating
a Euclidean version of the soft function is critical. Since the soft function in fact is a cross section and hence real and positive,
it satisfies the necessary condition for a Monte Carlo simulation. In this subsection, we present an approach to calculate it
in heavy-quark effective theory (HQET)~\cite{Ji:2019sxk}. There is also another method proposed to extract the reduced soft function $S_r$ from a light-meson form factor~\cite{Ji:2019sxk}, where many subtleties of HQET can be avoided. The first lattice calculation of the reduced soft function based on the light-meson formalism has been performed in the recent work~\cite{Zhang:2020dbb}.

%------------------------------------------------------------------------------
Due to the different space-time pictures of the DY and SIDIS processes, the soft functions for the two processes also differ from each other as shown in Fig.~\ref{Fig:soft1}. To define the soft function, one also needs to specify a time-ordering prescription. Since it is a cross section, it involves a time order and an anti-time order (or cut diagrams). However, in the light-cone limit, the time order does not matter.
What really matters is the rapidity regularization scheme. It has been proven for the $\delta$ regulator in ~\cite{Vladimirov:2017ksc} that the time ordering is not quite relevant up to overall phase factors, and the soft functions for the two processes are equal. The method therein can be modified to apply to the off-light-cone scheme too. Therefore, our first step is to convert the cut-diagrams into Feynman diagrams by imposing just
the single time order. In this way, the soft function can be viewed as a scattering amplitude.

In the off-light-cone scheme, there are further complications caused by the space-like or
time-like choices for off-light-cone vectors. Fortunately, one can show that in the light-cone limit, the space-like and time-like choices
are equivalent up to overall phase factors~\cite{future}. Thus we will use the notation $S(b_\perp,\mu,Y,Y')$ to denote a generic off-light-cone soft function that satisfies our demands.

With these in mind, we show that the off-light-cone soft function $S(b_\perp,\mu,Y,Y')$ is equivalent to an equal-time form factor of fast-moving color sources and can be formulated on the Euclidean lattice. From the matching formula Eq.~(\ref{eq:quasiTMDmatch}) in the last subsection, once the off-light-cone soft function is known, we can combine it with the lattice calculated quasi-TMDPDF to obtain the physical TMDPDF. Therefore, the cross section of DY processes in the low transverse-momentum region~\cite{Collins:1984kg} becomes predictable from first principles~\cite{Ji:2019sxk}.

To begin with, we define the scattering amplitude of a Wilson loop as shown in Fig.~\ref{fig:S}:
\begin{align}\label{eq:W_t}
	&W(t,t',b_\perp,Y,Y')\nonumber \\
	&=\frac{1}{N_c}{\rm Tr\,}\langle0|{\cal T} \left[{\cal W}^{\dagger}_{v'}(\vec{b}_\perp,t'){\cal W}_{v}(\vec{b}_\perp,t)\right]|0\rangle
\end{align}
where $|0\rangle$ is the QCD vacuum state and $N_c$ is number of colors and ${\rm Tr}$ is the color-trace.
Timelike four-vectors $v^\mu=\gamma(1,\beta,\vec 0_\perp)$ and $v'^\mu=\gamma'(1,-\beta',\vec 0_\perp)$ approach the lightcone as $\beta$ and $\beta' \to 1$. The rapidity $Y$ and the speed $\beta$ are related through $\beta = {\rm tanh} Y$, in terms of the light-cone vectors $p$ and $n$, the velocities read $v=\frac{e^Y}{\sqrt{2}}\left(\frac{p}{p^+}+e^{-2Y}p^+n \right)$ and $v'=\frac{e^Y}{\sqrt{2}}\left(e^{-2Y}\frac{p}{p^+}+ p^+ n \right)$. The ${\cal W}_{v}(\vec{b}_\perp,t)$ is a staple-shaped gauge-link along $v$ direction similar to those defined in Eqs.~(\ref{eq:staplep}) and (\ref{eq:staplez}). $t$ and $t'$ are the lengths of the $t$-components of the staples.
%
The single time-order prescription for $S$ allows physical interpretation as a chronological process.
\begin{figure}[htb]
	\centering
	\includegraphics[width=0.7\columnwidth]{images/tmd-S_fig13.pdf}
	\caption{The Wilson-loop ${\cal W}$ showing a pair of quark and antiquark scattering at $t=0$.}
	\label{fig:S}
\end{figure}
Similar to the quasi-TMDPDF, the Wilson-loop in Eq.~(\ref{eq:W_t}) contains pinch-pole singularities associated to time evolution of initial and final states at large $t$ and $t'$.
%
Therefore, we need to subtract them out in Eq.~(\ref{eq:W_t}) with rectangular Wilson-loops~\cite{Collins:2008ht,Ji:2018hvs}.
%
This leads to an off-light-cone realization of the soft function:
\begin{align}\label{eq:S_t}
	& S(b_\perp,\mu,Y,Y')\nonumber \\
	&=\lim_{\substack{t\to\infty\\t'\to\infty}}\frac{W(t,t',b_\perp,\mu,Y,Y')}{\sqrt{Z(2t,b_\perp,\mu,Y)Z(2t',b_\perp,\mu,Y')}} \ ,
\end{align}
where $Z(2t,b_\perp,Y)$ is the vacuum expectation value of rectangular Wilson loop which is similar to $W$ by setting $v'=v$ and $t'=t$, i.e. $Z(2t,b_\perp,Y)=W(t,t,b_\perp,Y,-Y)$.
%
The factor $Z$ has a clear physical interpretation: It can be viewed as the wave function renormalization for incoming or outgoing color sources.
%
After the subtraction through $Z$, the only remaining UV divergences for $S(b_\perp,\mu,Y,Y')$ are the cusp divergences with hyperbolic angle $Y+Y'$.

We should mention that a more common definition of the soft function $S_{\rm DY}(b_\perp,\mu,Y,Y')$ for the DY process was proposed in~\cite{Collins:2011ca,Collins:2011zzd}. The space-like vectors $u^\mu=\gamma(\beta,1,0,0)$ and $u'^\mu=\gamma'(-\beta',1,0,0)$ were chosen instead of time-like $v $ and $v'$ to define the soft function for the DY process. This soft function has already been defined in the last subsection in Eq.~(\ref{eq:S_DY}). $u$ and $u'$ are equal to $p_Y$, $n_{Y'}$ up to overall normalization factors.

While $S$ and $S_{\rm DY}$ are defined differently, we can show that
\begin{align}\label{eq:S_t=S}
	S(b_\perp,\mu,Y,Y')=S_{\rm DY}(b_\perp,\mu,Y,Y')
\end{align}
using analyticity property~\cite{future}. Here we focus on $S$ in Eq.~(\ref{eq:S_t}), which has a simple Euclidean realization.

After defining the soft function $S$, we now show that it is equal to a form factor. In HQET, the propagator of a color source is equivalent to a gauge-link along its moving direction.
%
Thus $W(t,t',b_\perp,\mu,Y,Y')$ can be expressed by fields in HQET with the Lagrangian
\begin{align}
	{\cal L_{\rm HQET}}=Q_v^{\dagger}(x)(iv \cdot D)Q_v(x)+ \bar Q_v^{\dagger}(x)(iv \cdot D)\bar Q_v(x) \ ,
\end{align}
where $Q_v$ and $\bar Q_v$ are quark and anti-quark in the fundamental and anti-fundamental representations, respectively; $v^\mu=\gamma(1,\beta,\vec 0_\perp)$ is the four velocity; $D$ is the covariant derivative.
%
Note that quarks in HQET can be viewed as color sources.
%
If the gluon soft function is considered, the heavy quarks should be in the adjoint representation.


In HQET, a color-singlet heavy-quark pair separated by $\vec b$ generates a heavy quark potential $V(|\vec b|)$ in the ground state, and the spectrum includes a gapped continuum above it.
%
The state can also have a residual momentum $\delta\vec P$, which is arbitrary due to reparameterization invariance~\cite{Luke:1992cs,Manohar:2000dt}, and for simplicity we always consider $\delta\vec P=0$.
%
When the sources move with a velocity $v$, the ground state can be labeled by
$|\overline Q Q,\vec{b},\delta \vec P\rangle_v$, where the residue momentum $\delta\vec P=\vec{P}_{\rm total}-2m_Q\gamma\vec\beta$ is the difference between the total momentum $\vec{P}_{\rm total}$ and the kinetic momentum of the heavy-quarks.
%
The residual energy of the state is $E=\gamma^{-1}V(|\vec b_{\perp}|)+\vec{\beta}\cdot \delta\vec P$.

Consider a process with incoming and outgoing states being heavy-quark pairs separated by ${\vec b}_\perp$ and at velocity $v$ and $v'$, respectively.
%
Such a state is created by the interpolating fields
\begin{align}
	{\cal O}_v(t,{\vec b}_\perp)=\int d^3\vec r \, Q_v^{\dagger}(t,\vec r\,) {\cal U}(\vec r,\vec r\,',t) \bar Q_v^{\dagger}(t,\vec r\,')  \ ,
\end{align}
where $\vec r\,'=\vec r+\vec b_\perp$ and ${\cal U}(\vec r,\vec r\,',t)$ is a gauge-link connecting $\vec r\,'$ to $\vec r$ at time $t$.
%
The heavy-quark pair created by ${\cal O}_v$ is forced to be at relative separation $\vec{b}_\perp$ and to have vanishing residual momentum $\delta \vec P=0$.
%
Between the incoming and outgoing states, a product of two local equal-time operators
\begin{align}
	J(v,v',\vec{b}_\perp)= \bar Q^\dagger_{v'}(\vec{b}_\perp)\bar Q_v(\vec{b}_\perp)Q_{v'}^\dagger(0)Q_v(0)
\end{align}
is inserted at $t=0$.
%
Then $W$ can be expressed in terms of HQET propagators which are gauge-links in the $v,v'$ directions.
%
After integrating out the heavy-quark fields, we obtain up to an overall volume factor~\cite{Ji:2019sxk}
\begin{align}\label{eq:W_HQ}
	&W(t,t',b_\perp,\mu,Y,Y')  \\
	&=\frac{1}{N_c} \langle0|{\cal O}^{\dagger}_{v'}(t',{\vec b}_\perp) J(v,v',\vec{b}_\perp) {\cal O}_v(-t,\vec{b}_\perp)|0\rangle \nonumber\\
	&\xrightarrow[\substack{t\to\infty\\t'\to\infty}]{} \frac{1}{N_c}\Phi^\dagger(b_\perp,\mu) S(b_\perp,\mu,Y,Y') \Phi(b_\perp,\mu)e^{-iE't'-iEt}  \ , \nn
\end{align}
where
\begin{align}\label{eq:S_HQ}
	\Phi(b_\perp,\mu)&=\lim_{T\to\infty}{}_v\langle \overline QQ,\vec{b}_\perp|{\cal O}_v(T,\vec{b}_\perp)|0\rangle\,,\\
	S(b_\perp,\mu,Y,Y')&={{}_{v'}}\langle\overline QQ,\vec{b}_\perp|J(v,v',\vec{b}_\perp)|\overline QQ,\vec{b}_\perp\rangle_v \ . \nn
\end{align}
%
In the last line of Eq.~(\ref{eq:W_HQ}), we have inserted a complete set of heavy-quark pair states before and after $J$.
%
At large $t$ and $t'$, the contribution from the continuum spectrum is damped out due to the Riemann-Lebesgue lemma~\cite{zuazo2001fourier}, while the contribution from $|\overline QQ,\vec{b}_\perp,\delta \vec{P}=0\rangle_v$ with residual energy $E=\gamma^{-1}V(|\vec b_\perp|)$ survives.
%
As a result we obtain Eqs.~(\ref{eq:W_HQ})---(\ref{eq:S_HQ}), where we have omitted the state label $\delta \vec{P}=0$ for simplicity.
%
Alternatively, we can also give $t$ and $t'$ a small negative imaginary part, which is consistent with the time order, to damp out all states except $|\overline QQ,\vec{b}_\perp\rangle_v$ at large $t$ and $t'$.
%
Note that $\Phi(\vec{b}_\perp,\mu)$ is independent of $Y$ because it is boost invariant.

Similarly, $Z$ can also be formulated in HQET as
\begin{align}\label{eq:Z_HQ}
	Z(2t,b_\perp,Y)&=\frac{1}{N_c}\langle0|{\cal O}^\dagger_{v'}(t,{\vec b}_\perp){\cal O}_v(-t,\vec{b}_\perp)|0\rangle \nonumber\\
	&\xrightarrow[\substack{t\to\infty}]{}\frac{1}{N_c}\Phi^\dagger({\vec b}_\perp,\mu)\Phi({\vec b}_\perp,\mu)e^{-2iEt}\,,
\end{align}
whose $t$-component has length $2t$.
%
The $Y$ dependence of $Z$ is implicit in the energy $E$.
%
Combining Eqs.~(\ref{eq:W_HQ}) and (\ref{eq:Z_HQ}), we obtain $S$ defined in Eq.~(\ref{eq:S_t}). We emphasize that Eq.~(\ref{eq:S_t}) can be seen as a LSZ reduction formula, in which we amputate the external heavy-quark pair states $|\overline QQ,{\vec b}_\perp\rangle_v$.

Being an equal-time observable, $S(b_\perp,\mu,Y,Y')$ can be straightforwardly realized in Euclidean time as:
\begin{align}\label{eq:S_lattice}
	&S(b_\perp,\mu,Y,Y') \\
	&=\lim_{\substack{T\to\infty\\T'\to\infty}}\frac{W_E(T,T',b_\perp,\mu,Y,Y')}{\sqrt{Z_E(2T,b_\perp,\mu,Y)Z_E(2T',b_\perp,\mu,Y')}}  \ , \nn
\end{align}
where the subscript $E$ indicates the quantity is defined in Euclidean time, with corresponding
variables $T$ and $T'$. Due to boost invariance, the factor $Z_E(T,b_\perp,\mu,Y)$ relates to the rectangular Wilson-loop defined in Eq.~(\ref{eq:Z_E}) along the $n_z$ direction through the relation $Z_E(2T,b_\perp,\mu,Y)=Z_E(2\gamma^{-1}T,b_\perp,0)$.
%
The relevant matrix elements are now calculated by a lattice version of HQET with the Lagrangian~\cite{Aglietti:1993hf,Hashimoto:1995in,Horgan:2009ti}
\begin{align}\label{eq:L_HQ}
	&{\cal L}_{{\rm HQET}}^E  \\
	&=Q_v^\dagger(x) (i\tilde v\cdot D_E)Q_v(x)+\bar Q_v^\dagger(x) (i\tilde v\cdot D_E)\bar Q_v(x) \ , \nn
\end{align}
where the subscript $E$ denotes the Euclidean space, $i\tilde v\cdot D_E=\gamma(D^\tau-i\beta )D^z$ with $\tilde v^\mu=\gamma(-i,-\beta,\vec 0_\perp)$. We have explicitly verified Eq.~(\ref{eq:S_lattice}) in Euclidean perturbation theory to the one-loop order.

The soft function cannot be calculated on the lattice by simply replacing the Minkowskian gauge-links in Eq.~(\ref{eq:W_t}) with a finite number of Euclidean gauge-links. Through HQET, we find a time-independent formulation of the soft function, which opens up the possibility of direct lattice calculations.

\subsection{Light-Front Wave-Function Amplitudes And Soft Function from Meson Form Factor}
\label{sec:others}
%------------------------------------------------------------------------------

Light-front quantization (LFQ) or formalism is a natural language for parton physics
in which partons are made manifest at all stages of calculations. It
favors a Hamiltonian approach to QCD like for a non-relativistic quantum
mechanical system, i.e., to diagonalize the Hamiltonian
\begin{equation}
	\hat P^-|\Psi_n\rangle = \frac{M^2_n+{\vec P_\perp}^2}{2P^+}|\Psi_n\rangle \ ,
\end{equation}
to obtain wave functions for the QCD bound states~\cite{Brodsky:1997de}.
The light-front wave functions (LFWFs) thus obtained can, in principle, be used to calculate
all the partonic densities and correlations functions. Moreover, like in
condensed matter systems, knowing quantum many-body wave-functions
allows one to understand interesting aspects of quantum coherence and
entanglement, as well as the fundamental nature of quantum systems.
Therefore, a practical realization of light-front quantization program clearly would be a big step
forward in understanding the fundamental structure of the proton.


However, from a field theory point of view, wave functions are not the most natural
objects to consider due to the non-trivial vacuum, UV divergences as well as the requirement of Lorentz symmetry, according to which the space and time should be treated on equal footing. The proton or other hadrons are
excitations of the QCD vacuum which by itself is very complicated
because of the well-known phenomena of chiral symmetry breaking and color confinement.
To build a proton on top of this vacuum, one naturally has a question of what part of the
wave-function reflects the property of the proton and what reflects the vacuum:
It is the difference that yields the properties
of the proton that are experimentally measurable. There is no clean way to make
this separation unless one builds the proton out of elementary excitations or quasi-particles that
do not exist in the vacuum, as often done in condensed matter systems.

The partons in the IMF avoids the above problems to a certain extent.
In fact, due to the kinematic effects, in the IMF all partons in
the vacuum have longitudinal momentum $k^+=0$, and to some degree of accuracy,
the proton is made of partons with $k^+\ne 0$. This natural separation of
degrees of freedom (DOF) is particularly welcome, making
a wave-function description of the proton more natural and interesting in IMF
than in any other frame.

To implement the above DOF separation, one possibility is to
assume triviality of the light-front vacuum. The question that to what extent this holds
has been continuously debated over the years.
One knows {\it a priori} that in relativistic QFT, the vacuum state is
boost invariant and frame-independent. In fact, it was proven in~\cite{Nakanishi:1976yx,Nakanishi:1976vf}
that not only the vacuum can not be trivial, even the Green's functions of the full theory cannot
pose generic meaningful restrictions to the null-planes $\xi^+=c$.  In fact, the vacuum
zero modes do contain non-trivial dynamics and contribute to the properties of the proton~\cite{Ji:2020baz}.
Nevertheless, one can adopt an effective theory point of view to simply cut off the zero-modes
and relegate their physics to renormalization constants. In some simple cases, these zero-modes can be
treated explicitly~\cite{Heinzl:1991myy,Yamawaki:1998cy}.

By imposing an IR cut-off on the $k^+\ge \epsilon$ in the effective Hilbert space,
all physics below $k^+=\epsilon$ are taken into account through
renormalization constants. We then obtain an effective LF theory with trivial
vacuum,
\begin{equation}
	a_{k\lambda}|0\rangle = b_{p\sigma}|0\rangle = d_{p\sigma}|0\rangle=0 \ .
\end{equation}
where $|0\rangle $ is the vacuum of LFQ.
Therefore, the proton can be expanded in terms of
the superposition of Fock states in the LF gauge $A^+=0$~\cite{Brodsky:1997de},
\begin{align}
	|P\rangle=\sum_{n=1}^{\infty} \int d\Gamma_n \psi^0_n(x_i,\vec{k}_{i\perp})\prod a^{\dagger}_i(x_i, \vec{k}_{i\perp})|0\rangle \ .
\end{align}
where $a^\dagger$ are generic quarks and gluon quanta on the light-front, the phase-space integral reads $d\Gamma_n=\prod \frac{dk^+d^2k_\perp}{2k^+(2\pi)^3}$. The $\psi_n(x_i,\vec{k}_{i\perp})$
are LFWF amplitudes or simply WF amplitudes, where $x_i$ to denote the set of momentum fractions from $x_1$ to $x_n$.
They are a complete set of non-perturbative quantities which describe the partonic
landscape of the proton. The above amplitudes can in principle be calculated through Hamiltonian
diagonalization. However, as explained in Sec.~\ref{sec:lamet-frame}, a direct systematic
solution in LFQ is impractical.

LaMET offers an alternate route to calculate these WF amplitudes. Thanks to the triviality
of the vacuum after the truncation $k^+\ge \epsilon$, they can then be written in terms of the
invariant matrix elements by inverting the above expansion,
\begin{align}
	\psi^0_n(x_i,\vec{k}_{i\perp})=\langle 0|\prod a_i(x_i,\vec{k}_{i\perp})|P\rangle \ .
\end{align}
After properly restoring gauge-invariance and imposing regularizations, they
become the matrix elements of light-cone correlators, the same type as those in the
TMDPDFs. Therefore, the LaMET method applies to them, which allows one to effectively
obtain the results of light-front quantization through instant quantization in a large momentum frame.

To realize the goal, the LFWF amplitudes also need a rapidity renormalization, as in the case of TMDPDFs.
In this section, we explain how the reduced soft function $S_r$ can be obtained by combining the LFWF amplitudes
and a special light-meson form factor, instead of as the form factor in HQET discussed in the previous section. A lattice calculation based on the light-meson framework has been performed in~\cite{Zhang:2020dbb}.

Let us consider the following form factor of a pseudoscalar light-meson state with constituents $\overline\psi\eta$,
\begin{align}
	F(b_\perp,P,P',\mu)=\langle P'|\overline\eta(\vec b_\perp)\Gamma'\eta(\vec b_\perp) \overline\psi(0)\Gamma\psi(0)|P\rangle
\end{align}
where $\psi$ and $\eta$ are light quark fields of different flavors; $P^\mu=(P^t,0,0,P^z)$ and $P'^\mu=(P^t,0,0,-P^z)$ are two large momenta which approach two opposite light-like directions in the limit $P^z\to\infty$;
$\Gamma$ and $\Gamma'$ are Dirac gamma matrices, which can be chosen as $\Gamma=\Gamma'=1$, $\gamma_5$ or
$\gamma_\perp$ and $\gamma_\perp\gamma_5$, so that the quark fields have leading components
on the respective light-cones.

At large momentum, the form factor factorizes through TMD factorization into LFWF amplitudes. To motivate the factorization,
we need to consider the leading region of IR divergences in a similar way for SIDIS and Drell-Yan~\cite{Ji:2004wu,Collins:2011ca},
and the result is shown in Fig.~\ref{fig:reduced_form}.
\begin{figure}
	\includegraphics[width=0.7\columnwidth,angle=90]{images/wf-reduced_form_fig14.pdf}
	\caption{\label{fig:reduced_form} The reduced diagram for the large-momentum form factor $F$ of
		a meson. Two $H$ denote the two hard cores separated in space by $\vec{b}_\perp$, $C$ are collinear sub-diagrams
		and $S$ denotes the soft sub-diagram.}
\end{figure}
There are two collinear sub-diagrams responsible for collinear modes in $+$ and $-$ directions, and a soft sub-diagram responsible for soft
contributions. Besides, there are two IR-free hard cores localized around $(0,0,0,0)$ and $(0,\vec{b}_\perp,0)$. In the covariant gauge, there are arbitrary numbers of longitudinally-polarized collinear and soft gluons that can connect to the hard and collinear sub-diagrams.
Based on the region decomposition, we now follow the standard procedure to make factorization into
LF quantities~\cite{Collins:2011ca}.

We first factorize the soft divergences. This can be done with the soft function $S(b_\perp,\mu,\delta^+,\delta^-)$. It re-sums the soft gluon radiations from fast-moving color-charges. Intuitively, soft gluons have no impact on the velocity of the fast-moving color charged partons, and the propagators of partons eikonalize to straight gauge links along their moving trajectory.

We then factorize the collinear divergences. For the incoming direction, the collinear divergences is captured by the LFWF amplitude for the incoming parton $\psi_{\bar q q}(x,b_\perp,\mu,\delta^{'-})$ defined with future-pointing gauge-links.
\begin{align}
	&\psi_{\bar q q}(x,b_\perp,\mu,\delta^{'-})=\int\frac{d\lambda}{4\pi}e^{-ix\lambda}\\
	&\langle 0|\bar \psi(\lambda n/2 +\vec{b}_\perp)\gamma^+{\cal W}_n(\lambda n/2+\vec{b}_\perp)|_{\delta^{'-}}\psi(-\lambda n/2)|P\rangle \ ,\nn
\end{align}
where the staple-shaped gauge-link $W_n$ is defined similar to that in Eq.~(\ref{eq:staplen}), the only exception being the gauge-links $W_n$ should point to $+\infty$ instead of  $-\infty$.

However, the naive LFWF amplitude contains soft divergences as well, to avoid double-counting, we must subtract out the soft contribution from the bared collinear WF amplitude with the soft function $S(b_\perp,\mu,\delta^+,\delta^{'-})$. This leads to the collinear function for the incoming direction: $\psi_{\bar q q}(x,b_\perp,\mu,\delta^{'-})/S(b_\perp,\mu,\delta^+,\delta^{'-})$. Similarly, for the out-going direction one obtains the collinear function $\psi^{\dagger}(x',b_\perp,\mu,\delta^{'+})/S(b_\perp,\mu,\delta^{'+},\delta^{-})$.

Here we briefly comment on the choices for the gauge-link directions in the soft functions and the WF amplitudes. Naively, the gauge-links along the $p$ direction have to be past-pointing. However, similar to the arguments in~\cite{Collins:2004nx} for the SIDIS process, based on the space-time picture of collinear divergences, one can chose future pointing gauge-links along $p$ direction as well. With all the gauge-links being future pointing, the soft function equals to $S^-$ which is manifestly real, and the WF amplitudes for the incoming and outgoing hadrons are in complex conjugation to each other.

Besides the collinear and soft functions, we still need the hard core $H_F(Q^2,\bar Q^2,\mu^2)$ where $Q^2=xx'P\cdot P'$, $\bar Q^2=\bar x\bar x' P\cdot P'$ and an integral over the momentum fractions $x$,$x'$ is assumed. Taking together, we have the TMD factorization of the form factor into hard, collinear and soft functions:
\begin{align}\label{eq:form_fac_bare}
	&F(b_\perp,P,P',\mu)=\int dx dx' H_F(Q^2,\bar Q^2,\mu^2) \\ &\times \left[\frac{\psi_{\bar q q}^{\dagger}(x',b_\perp,\mu,\delta^{'+})}{S(b_\perp,\mu,\delta^{'+},\delta^-)}\right]\left[\frac{\psi_{\bar q q}(x,b_\perp,\mu,\delta^{'-})}{S(b_\perp,\mu,\delta^+,\delta^{'-})}\right]\nonumber \\
	&\times S(b_\perp,\mu,\delta^+,\delta^-) \ . \nn
\end{align}
All the rapidity regulators in all the WF amplitudes and the soft functions are cancelled.

Let us consider a one-loop example. The incoming hadron state consists of a free quark with momentum $x_0P^+$ and a free anti-quark with momentum $\bar x_0 P^+$. Similarly the outgoing state consists of a pair of free quark and anti-quark with momentum $x'_0P^{'-}$, $\bar x'_0 P^{'-}$, respectively. The spin projection operator for the incoming state is proportional to $\gamma^5 \gamma^-$ and for the out-going state is proportional to $\gamma^5 \gamma^+$. The tree level form factor is normalized to 1. At one-loop level, the pseudo-scalar form factor with vector currents $\Gamma=\gamma^{\mu}$, $\Gamma'=\gamma_{\mu}$ where a summation over $\mu$ is assumed reads:
\begin{align}
	F(b_\perp,P,P',\mu)=1+\frac{\alpha_s C_F}{2\pi}F^{(1)}(b_\perp,Q^2,\bar Q^2,\mu^2)\ ,
\end{align}
where $Q^2=2x_0x_0'P^+P^{'-}$, $\bar Q^2=2\bar x_0\bar x_0'P^+P^{'-}$ and
\begin{align}
	&F^{(1)}(b_\perp,Q^2,\bar Q^2,\mu^2) \\
	&=-7+\left(-\frac{1}{2}\ln^2 b^2_\perp Q^2+\frac{3}{2}\ln b^2_\perp Q^2+\left(Q \rightarrow \bar Q\right)\right) \ . \nn
\end{align}
This result can be obtained from the one-loop DY structure function~\cite{DAlesio:2014mrz} using the substitution $\ln^2(-Q^2b^2_\perp)\rightarrow \frac{1}{2}\ln^2 Q^2b^2_\perp+\ln^2 \bar Q^2b^2_\perp$ and $\ln(-Q^2b^2_\perp)\rightarrow \frac{1}{2}\ln Q^2b^2_\perp+\ln \bar Q^2b^2_\perp$. Similar to the TMD factorization for SIDIS and DY process, one should also notice that the hard kernel $H_F(Q^2,\bar Q^2,\mu^2)$ can be obtained from that of the space-like Sudakov form factor:
\begin{align}
	H_F(Q^2,\bar Q^2,\mu^2)=H^{\rm sud}(-Q^2)H^{\rm sud}(-\bar Q^2) \ ,
\end{align}
where $H^{\rm sud}(-Q^2)$ is given in~\cite{Collins:2017oxh}. At one-loop level, we then obtain:
\begin{align}
	& H_F(Q^2,\bar Q^2,\mu^2)=1\nonumber \\
	&+\frac{\alpha_s}{4\pi}\left(-16+\frac{\pi^2}{3}+3L_Q+3L_{\bar Q}-L_Q^2-L_{\bar Q}^2\right)  \ ,
\end{align}
where $L_Q=\ln \frac{Q^2}{\mu^2}$ and $L_{\bar Q}=\frac{\bar Q^2}{\mu^2}$.

Now we construct the Euclidean version of the factorization in terms of the quasi-WF amplitudes,
the reduced soft function, and hard contribution. The quasi-WF amplitudes are defined in a way similar to Eq.~(\ref{eq:quasi_TMD}):
\begin{align}
	&\tilde \psi_{\bar q q}(x,b_\perp,\mu,\zeta_z)=\int \frac{d\lambda}{4\pi}e^{ix\lambda}\nonumber\\
	&\frac{\langle 0| \bar \psi\big(\frac{\lambda n_z }{2}\!+\!\vec{b}_\perp\big)\gamma^z{\cal W}_{z}(\frac{\lambda n_z}{2}\!+\!\vec{b}_\perp;L)\psi\big(\!-\!\frac{\lambda n_z}{2}\big) |P\rangle}{Z_E(2L,b_\perp,\mu)} \ ,
\end{align}
in which the staple-shaped gauge-link ${\cal W}_z$ is defined in Eq.~(\ref{eq:staplez}). The gauge-links should point to $+z$ direction in accordance to the $+\infty$ choice on the light-cone side.

The factorization to the LFWF amplitude follows a similar reasoning to that of the quasi-TMDPDFs presented in previous sections.  Alternatively, we can factorize it using quantities defined in on-light-cone rapidity scheme,
\begin{align}\label{eq:quasifac_bare}
	\tilde \psi_{\bar q q}(x,b_\perp,\mu,\zeta_z)&=H^+_1\left(\zeta_z/\mu^2,\bar \zeta_z/\mu^2\right)\\ &\times \left[\frac{\psi_{\bar q q}(x,b_\perp,\mu,\delta^-)}{S(b_\perp,\mu,\delta^+,\delta^-)}\right]S(b_\perp,\mu,\delta^+) \ . \nn
\end{align}
This factorization is the result of applying a similar leading-region analysis to the quasi-WF amplitude. The $\psi_{\bar q q}(x,b_\perp,\mu,\delta^-)/S(b_\perp,\mu,\delta^+,\delta^-)$ re-sums all the collinear divergences, while the soft function $S(b_\perp,\mu,\delta^+)$ contains an off-light-cone direction along ${n_z}$. It re-sums the soft divergences of the quasi-WF amplitude. The soft functions $S(b_\perp,\mu,\delta^+,\delta^-)$ and $S(b_\perp,\mu,\delta^+)$ subtract away the regulator dependencies introduced in the bare LFWF amplitude. The overall combination in the right-hand side of Eq.~(\ref{eq:quasifac_bare}) is rapidity-scheme independent. Similar to the case of the form factor, we can chose all the gauge-links along the incoming collinear direction to be future-pointing.

Combining together Eqs.~(\ref{eq:form_fac_bare}) and (\ref{eq:quasifac_bare}) and using the relation $\zeta \zeta'= \zeta_z \zeta'_z$,
one obtains the form factor factorization,
\begin{align}
	& F(b_\perp,P,P',\mu)\\
	&= \int dx dx' H(x,x')\tilde \psi_{\bar q q}^{\dagger}(x',b_\perp)\tilde \psi_{\bar q q}(x,b_\perp)S_{r}(b_\perp,\mu)\ ,\nn
\end{align}
where we have only kept the $x,b_\perp$ dependencies of the WF amplitudes with other variables being omitted
, and the hard kernel $H$ is given by:
\begin{align}
	&H(x,x')=H(\zeta_z,\zeta'_z,\bar \zeta_z,\bar \zeta'_z,\mu^2)\nonumber \\
	&\qquad=\frac{H_F(Q^2,\bar Q^2,\mu^2)}{H^+_1\left(\zeta_z/\mu^2,\bar \zeta_z/\mu^2\right)H^+_1\left(\zeta'_z/\mu^2,\bar \zeta'_z/\mu^2\right)} \ ,
\end{align}
where $Q^2=\sqrt{\zeta_z\zeta'_z}$ and $\bar Q^2=\sqrt{\bar \zeta_z\bar\zeta'_z}$. And the reduced soft function is
\begin{align}
	S_{r}(b_\perp,\mu)=\lim_{\delta^+,\delta^-\rightarrow 0}\frac{S(b_\perp,\mu,\delta^+,\delta^-)}{S(b_\perp,\mu,\delta^+)S(b_\perp,\mu,\delta^-)} \ .
\end{align}
It can be shown based on property of off-light-cone soft functions that $S_r$ defined here agrees with the one defined in Eq.~(\ref{eq:reduced}).

Therefore, with non-perturbative quantities $F$ and $\psi^+$, we obtain the reduced soft function,
\begin{equation}
	S_{r}(b_\perp,\mu)=\frac{F(b_\perp,P,P',\mu)}{\int dx dx' H(x,x')\tilde \psi_{\bar q q}^{\dagger}(x',b_\perp)\tilde \psi_{\bar q q}(x,b_\perp)} \ ,
	\label{eq:form_qfac}
\end{equation}
where $H$ can be obtained perturbatively.

Based on the one-loop results for the form factor, the quasi-WF amplitudes and the reduced soft function, the one-loop matching kernel for the vector current can be extracted as:
\begin{align}
	&H(\zeta_z,\zeta'_z,\bar \zeta_z,\bar \zeta'_z,\mu^2)=1+ \frac{\alpha_sC_F}{2}i\ln \frac{\sqrt{\zeta_z\bar\zeta_z}}{\sqrt{\zeta'_z\bar\zeta'_z}}  \\
	&+\frac{\alpha_sC_F}{4\pi}\left(-8+\ln^2\frac{\sqrt{\zeta_z}}{\sqrt{\zeta'_z}}+\ln \frac{\sqrt{\zeta_z\zeta'_z}}{\mu^2}+(\zeta \rightarrow \bar \zeta)\right)\nonumber \ ,
\end{align}
and the renormalization group equation for $H$ reads:
\begin{align}
	\mu^2\frac{d}{d\mu^2}\ln H(\zeta_z,\zeta'_z,\bar \zeta_z,\bar \zeta'_z,\mu^2)=-2\gamma_F(\alpha_s)-\Gamma_{S}(\alpha_s) \ ,
\end{align}
where $\gamma_F$ and $\Gamma_S$ have been defined before.

Here we briefly comment on the end-point behavior. As $x\sim 0$, the hard kernel diverges logarithmically near the end point as $ 1+\alpha_s \ln^2 x$, but the quasi-WF amplitudes approach zero at large or small $x$ linearly, thus the end point regions behave as $x\ln^2 x$, which is free from those problems for the $k_T$ factorization for electromagnetic form factor~\cite{Li:1992nu}. Moreover, we can fix the $z$-component momentum transfer at each of the vertices to be $P^z$, which indicating that $x+x'=1$. In this case the end-point behavior is improved to $x^2\ln^2 x$.

%------------------------------------------------------------------------------
